\chapitre{Vectorisation des colonnes sur processeur}
La boucle interne de NvPD calcule une ligne à partir de la table des dernières occurrences. Le branchement sur \(k=0\), \(k=j\) ou \(0<k<j\) peut être remplacé par le minimum de l'algorithme 1. Quand \(k=j\), le second terme vaut \(P_{j-1}\), le coût d'une correspondance. Le cas sans correspondance suit la borne établie pour l'algorithme 1. Une instruction SIMD peut alors traiter plusieurs colonnes qui choisissent des prédécesseurs différents.
\chart{03_variantes_cpu.png}{Les rapports de temps comparent les noyaux sur des paires identiques à largeur de ligne et nombre de cœurs fixés.}{fig:simd}{bench/profile_cpu_driver.cpp}
Les essais de la figure~\ref{fig:simd} utilisent GCC~11.5 avec \texttt{-O3 -march=native} sur EPYC~9654. La médiane des rapports \(T_\mathrm{initial}/T_\mathrm{SIMD}\) porte sur trois répétitions appariées. Le compilateur vectorise la boucle sur 32 octets, soit huit entiers par instruction vectorielle. À un cœur, le gain atteint 5,85 sur la forme \(256\times10^6\). Dans le panneau b) de la figure~\ref{fig:million}, la boucle initiale et la boucle SIMD utilisent les mêmes huit cœurs et les mêmes options de compilation. À 2\,\%, le temps passe de 324,41 à 53,04\,s pour \(10^6\times1\,000\,033\) cases utiles, soit 18,85\,GCUPS pour le noyau SIMD. Pour les trois taux, les gains vont de 6,12 à 6,29 avec un seul passage de chaque variante par paire.
\chart{05_cout_instructions.png}{Les compteurs de l'EPYC rapportent les instructions retirées et les erreurs de prédiction au nombre de cases des trois matrices.}{fig:instructions}{bench/profile_cpu.py}
La figure~\ref{fig:instructions} rapporte les compteurs au travail dense \(W=mn\). Sur un seul cœur et \(20\,000^2\) cases, la reformulation ramène les instructions retirées de 21,76 à 2,13 par case. Les cycles comprennent aussi les dépendances de données et les attentes mémoire.
\synthese{Sur le carré de 20\,000 bases, la reformulation accélère le calcul à un cœur d'un facteur 2,32.}{La lecture indirecte en \(k-1\) conduit à distinguer l'effet des instructions de celui du cache.}

\chapitre{Empreinte de l’index et localité sur processeur}
La taille de l'index détermine la part du jeu de données qui tient en cache. Pour \(q\) lettres distinctes et \(n\) colonnes, l'index et les deux lignes du noyau de coût seul occupent
\[
 M_{\mathrm{NvPD}}=4(q+2)(n+1)\quad [\text{octets}].
\]
Les deux lignes sont contiguës. Le dernier prédécesseur varie avec la lettre. Le prototype compact code l'écart dans un octet lorsque sa valeur et les sentinelles tiennent dans cette représentation. La conversion garde provisoirement l'index entier. Chaque nœud CPU réunit deux EPYC 9654 de 96 cœurs. Huit cœurs partagent 32\,Mio de cache L3. Chaque cœur dispose de 32\,Kio de L1 de données et de 1\,Mio de L2.
\chart{04_memoire_cpu.png}{Le pic de mémoire et les défauts L1 de lecture sont comparés aux mêmes dimensions de matrice sur huit cœurs.}{fig:mem}{bench/profile_cpu.py}
Sur la forme \(64\times4\,194\,304\), la figure~\ref{fig:mem} distingue le maximum d'espace de travail, qui passe de 96 à 80\,Mio, du RSS\footnote{Le maximum de mémoire physique résidente comprend le programme, ses entrées et les allocations des bibliothèques.}. Sur le carré à un cœur, le nombre de défauts L1 pour mille cases varie de 188,97 à 189,39 entre la boucle initiale et la boucle SIMD. Un défaut L1 peut être servi par le L2 ou le L3.
\synthese{Sur la forme la plus large, l'index compact ramène le RSS d'environ 101,5 à 85,0\,Mio.}{Le rapport des temps SIMD/compact va de 0,88 à 1,14 sur les neuf configurations, ce qui motive l'étude du placement des fils.}

\chapitre{Passage à l’échelle en mémoire partagée}
Le passage à l'échelle forte répartit un travail constant entre \(p\) cœurs. Le passage à l'échelle faible garde le nombre de cases par cœur constant~\cite{amdahl,gustafson}. Pour une matrice de travail \(W\), les indicateurs valent
\[
 S_p=\frac{T_1(W)}{T_p(W)},\qquad E_p=\frac{S_p}{p},\qquad
 E_p^{\mathrm{faible}}=\frac{T_1(W_1)}{T_p(pW_1)}.
\]
La série forte utilise \(20\,000^2\) cases. La série faible conserve 4\,096 lignes pour \(4\,096p\) colonnes. Elle fixe donc le travail à 16\,777\,216 cases par cœur et élargit les tableaux avec \(p\).
\chart{02_echelle_cpu.png}{Les courbes donnent l'accélération forte et l'efficacité faible de NvPD selon le nombre de cœurs physiques de l'EPYC.}{fig:echellecpu}{bench/profile_cpu.sbatch}
À taille fixe, huit cœurs prennent 27,93\,ms contre 31,28\,ms à seize sur la figure~\ref{fig:echellecpu}. Au dernier placement, les cœurs occupent deux groupes de cache et deux domaines NUMA.\footnote{La distance entre les cœurs et les contrôleurs mémoire dépend du domaine NUMA auquel chacun appartient.} Les synchronisations entre lignes et les accès mémoire interviennent dans le même temps mesuré.
\synthese{Huit cœurs atteignent 73,6\,\% d'efficacité forte sur la paire carrée.}{L'efficacité faible descend à 62,4\,\% à seize cœurs, avec des lignes seize fois plus larges qu'à un cœur.}
